\uuid{dqZ5}
\titre{Rang d'une famille de vecteurs}
\auteur{hanani ; mzari}
\organisation{exo7}
\licence{CC BY-NC-SA 4.0}
\datecreate{2020-01-20}
\niveau{L1}
\module{Algèbre}
\chapitre{Espace vectoriel}
\sousChapitre{Base}
\difficulte{2}
\theme{Base et dimension, Rang d'une famille de vecteurs}
\exo7qcmid{lille-2.espaces-vectoriels.029}

\contenu{
\begin{question}
\qclassification{106.04}
\qtags{motcle=Rang d'une famille de vecteurs}
Dans $\Rr^3$, on considère les vecteurs
$$u_1=(-1,1,2),\quad u_2=(0,1,1),\quad  u_3=(-1,0,1),\quad u_4=(0,2,1).$$
Quelles sont les assertions vraies ?
\begin{answers}
\good{Le rang de la famille $\{u_1,u_2\}$ est $2$}.
\bad{Le rang de la famille $\{u_1,u_2,u_3\}$ est $3$}.
\bad{Le rang de la famille $\{u_1,u_2,u_3,u_4\}$ est $4$}.
\good{Le rang de la famille $\{u_1,u_2,u_4\}$ est $3$}.
\end{answers}
\begin{explanations} Le rang d'une famille de vecteurs est la dimension du sous-espace vectoriel
engendré par ces vecteurs. Autrement dit, c'est le nombre maximum de vecteurs de cette famille qui sont linéairement indépendants. On vérifie que $u_1=u_2+u_3 $ et que $\{u_1,u_2,u_4\}$ est libre.
\end{explanations}
\end{question}
}
